\chapter{Attached files}
The submitted attachments are:
\begin{itemize}
    \item code.zip
    \item 3d\_models.zip
\end{itemize}

\section{code.zip}
The structure of \texttt{code.zip} is as follows:
\begin{verbatim}
code
|-- README.md
|-- electronic_design.pdf
|-- examples
|   |-- pcd_2026.05.16_13:41:33_495468-badalign-nocompensation.ply
|   |-- pcd_2026.05.16_13:54:25_494384-badalign-minus1scanangle-ceilingcorrect.ply
|   `-- pcd_2026.05.16_14:02:09_494400-best-minus1scanangle-plus1.5z.ply
|-- pc.py
|-- pcd_visualizer.py
|-- pcd_visualizer_outlier_removal.py
|-- pico_load
|   |-- main.py
|   |-- secrets.py
|   `-- stepper.py
`-- requirements.txt
\end{verbatim}
\begin{itemize}
    \item \texttt{README.md} contains instructions for setting up the Python \texttt{.venv} environment and running the project.
    \item \texttt{electronic\_design.pdf} is a wiring diagram created in Fritzing that shows the pin connections used in this project.
    \item \texttt{examples/} folder contains point clouds that were discussed in \cref{subsec:misalignment_corrections}.
    \item \texttt{pc.py} is the PC code responsible for communication with Pico, acquisition and parsing of \gls{lidar} data, point cloud reconstruction and management of the visualisation window showing the real-time reconstruction.
    \item \texttt{pcd\_visualizer.py} and \texttt{pcd\_visualizer\_outlier\_removal.py} are scripts used for visualising the stored point clouds
    \item \texttt{pico\_load/} folder contains code intended to be uploaded to Pico. \texttt{secrets.py} should be filled with the preferred Wi-Fi SSID and password used for connecting to the Pico access point.
    \item \texttt{requirements.txt} contains Python libraries needed for running the code on PC.
\end{itemize}

\section{3d\_models.zip}
\begin{verbatim}
3d_models
|-- Bottom.stl
|-- CCW Gear x6.stl
|-- DC housing.stl
|-- Lower Carrier.stl
|-- Motor Shaft Gear.stl
|-- Platform Mount.stl
|-- Platform.stl
|-- Ring Gear Lid.stl
|-- Ring Gear.stl
|-- Rod x6.stl
|-- Upper Carrier.stl
`-- model_attribution.txt
\end{verbatim}
The \texttt{3d\_models} contains all 3D models in \texttt{.stl} format used for building the platform. 
The suffix \texttt{x6} in some filenames indicates the corresponding number of copies required for the assembly.

The \texttt{DC housing.stl} is not author's work. The attribution link is provided in  \texttt{model\_attribution.txt}.

The 3D modelling was done in Onshape. The complete assembly can be accessed via \href{https://cad.onshape.com/documents/0d531e009de73b9ed4692d36/w/83d10e6192a92a79ef75f390/e/35ebef400ac87f1073d7eabc?renderMode=0&uiState=6a06f9cc76a0e4c598c450d3}{this link}.

\section{Project Repository}
The entire development of this thesis was managed using Git. The repository is hosted on the faculty GitLab server and can be accessed via \href{https://gitlab.fel.cvut.cz/novosill/bap-novosad/-/tree/4b53a596703105143e23482818c8093a5936ead5/}{this link}.